For the first conclusion, fix \(n\).  Since \(\delta_n>0\), specialize \(\text{lem:jms-ace-theorem-five-four}\) by replacing its positive separation parameter \(\delta\) with \(\delta_n\).  The resulting model is exactly \(\mathcal C_n(\delta_n)\), because no other member condition changes.  The quantities \(a_{1,n},b_{1,n},a_2\) are unchanged, whereas \(b_{2,n}\) becomes
\[
 \widetilde b_{2,n}
 =200\min\{1,C_{\theta}\}\delta_n
 \max\!\left\{\varepsilon_{1,n},\varepsilon_{2,n},
 (\gamma n)^{-1/2}(\psi_{\xi}+C_{\theta}\psi_{\eta})\right\}^{-1}.
\]
Thus the displayed local eligibility condition is precisely the source theorem's eligibility condition under this specialization.  Positivity of \(\varepsilon_{1,n},\varepsilon_{2,n}\), together with the explicit requirement that every logarithm and denominator be defined, supplies all boundary conditions of the cited result.  It follows that
\[
\begin{aligned}
 &\sup_{P\in\mathcal C_n(\delta_n)}
 Q_{P,1-\gamma}(|\widehat\theta_{\mathrm{ACE},r,n}-\theta_0(P)|)\\
 &\quad\le C_{\gamma}r!16^r\delta_n^{-1}
 \Bigl[\varepsilon_{1,n}^r\varepsilon_{2,n}
 +C_{\theta}\varepsilon_{1,n}^{r+1}\\
 &\hspace{35mm}
 +64(C_g+\psi_{\eta})^r
 \{r^2(C_g+\psi_{\eta})+\psi_{\xi}+C_{\theta}\psi_{\eta}\}
 (\gamma n)^{-1/2}\Bigr].
\end{aligned}
\]
This is \(U_{\mathrm{ACE},n}^{\mathrm{loc}}\).  Taking the smaller of this numerical upper bound and a separately valid DML upper bound is only an oracle comparison: it neither defines a sample-measurable selector nor supplies a lower bound.

For the second conclusion, the theorem's displayed \(L^s(P_X)\) premise, with \(s\ge r+1=4\), implies by Lyapunov's inequality
\[
 \|\bar g_n-g_0\|_{L^1(P_X)}
 \le\|\bar g_n-g_0\|_{L^s(P_X)}
 \le\varepsilon_{1,n}.
\]
All remaining displayed path assumptions coincide with those of \(\text{lem:gaussian-rademacher-l1-benchmark}\).  Applying that lemma yields the stated transform, cumulant, zero, annihilation identity, denominator lower bound, explicit clipped estimator, and both risk displays.  In particular, the proof establishes the path conclusion under the weaker direct \(L^1(P_X)\) condition; the frozen \(L^s(P_X)\) premise is sufficient but not consumed beyond its implication of \(L^1(P_X)\) control.  Since the calculation concerns only the explicit Gaussian--Rademacher path, it proves no converse and no uniform upper bound for the larger class defined solely by fourth-cumulant separation.  Hence neither conclusion determines the full minimax phase diagram.